\documentclass[12pt]{article}

\usepackage[a4paper,margin=1in]{geometry}
\usepackage{amsmath,amssymb,amsthm}
\usepackage{enumitem}
\usepackage{array}

\setlength{\parindent}{0pt}
\setlength{\parskip}{0.7em}

\newtheorem*{remark}{Remark}

\begin{document}

\begin{center}
    {\Large \textbf{Tutorial: Time Series Question 1}}\\
    \vspace{0.2cm}
    {\large ACF, PACF, AR Variance, and Stationarity}
\end{center}

\section*{How to Use This Tutorial}

This file explains one big question step by step.  Before each part, the relevant
theory is stated first.  The aim is not only to answer this exact question, but
also to show the method you can reuse for similar exam questions.

\section*{Question 1(a): ACF of an MA(1)-Type Process}

\subsection*{Theory to Use}

The autocovariance function is
\[
\gamma_k = \operatorname{Cov}(Y_t,Y_{t-k}).
\]

The autocorrelation function is
\[
\rho_k = \frac{\gamma_k}{\gamma_0}.
\]

For a moving average process of order 1,
\[
Y_t = e_t + \psi e_{t-1},
\]
where \(e_t\) is white noise with variance \(\sigma_e^2\), the autocovariances are
\[
\gamma_0 = (1+\psi^2)\sigma_e^2,
\]
\[
\gamma_1 = \psi \sigma_e^2,
\]
and
\[
\gamma_k=0,\qquad |k|\geq 2.
\]

So the ACF is
\[
\rho_1 = \frac{\psi}{1+\psi^2},
\]
and
\[
\rho_k=0,\qquad |k|\geq 2.
\]

\subsection*{Question}

Consider the model
\[
Y_t = e_t - \frac{1}{\theta}e_{t-1},
\]
where
\[
e_t \sim WN(0,\theta^2\sigma^2).
\]
Obtain the ACF of \(Y_t\).

\subsection*{Answer}

Compare the given model with the standard MA(1) form:
\[
Y_t = e_t + \psi e_{t-1}.
\]

Here,
\[
\psi = -\frac{1}{\theta}.
\]

Also,
\[
\operatorname{Var}(e_t)=\theta^2\sigma^2.
\]

Therefore,
\[
\gamma_0
=
\left(1+\psi^2\right)\operatorname{Var}(e_t).
\]

Substitute \(\psi=-1/\theta\):
\[
\gamma_0
=
\left(1+\frac{1}{\theta^2}\right)\theta^2\sigma^2.
\]

Simplifying,
\[
\gamma_0
=
(\theta^2+1)\sigma^2.
\]

Now find \(\gamma_1\):
\[
\gamma_1
=
\psi \operatorname{Var}(e_t).
\]

Thus,
\[
\gamma_1
=
\left(-\frac{1}{\theta}\right)\theta^2\sigma^2
=
-\theta\sigma^2.
\]

For an MA(1) process,
\[
\gamma_k=0,\qquad |k|\geq 2.
\]

Therefore,
\[
\rho_1
=
\frac{\gamma_1}{\gamma_0}
=
\frac{-\theta\sigma^2}{(\theta^2+1)\sigma^2}.
\]

So
\[
\boxed{\rho_1=-\frac{\theta}{\theta^2+1}}.
\]

Also,
\[
\boxed{\rho_0=1},
\]
and
\[
\boxed{\rho_k=0,\qquad |k|\geq 2}.
\]

Therefore, the ACF is
\[
\boxed{
\rho_k =
\begin{cases}
1, & k=0,\\[4pt]
-\dfrac{\theta}{\theta^2+1}, & |k|=1,\\[8pt]
0, & |k|\geq 2.
\end{cases}
}
\]

\subsection*{Simple Interpretation}

The process only contains \(e_t\) and \(e_{t-1}\).  So \(Y_t\) can only overlap
with \(Y_{t-1}\).  After lag 1, there is no shared white-noise term, so the ACF
becomes zero.

\section*{Question 1(b): Derive Lag 2 PACF}

\subsection*{Theory to Use}

The partial autocorrelation at lag 2 measures the direct relationship between
\(Y_t\) and \(Y_{t-2}\), after removing the effect of \(Y_{t-1}\).

To derive lag 2 PACF, use the Yule-Walker equations for an AR(2) approximation:
\[
Y_t = \phi_{21}Y_{t-1}+\phi_{22}Y_{t-2}+a_t.
\]

The Yule-Walker equations in terms of autocorrelations are
\[
\rho_1 = \phi_{21} + \phi_{22}\rho_1,
\]
\[
\rho_2 = \phi_{21}\rho_1 + \phi_{22}.
\]

The lag 2 PACF is
\[
\phi_{22}.
\]

\subsection*{Question}

Derive the lag 2 PACF for a time series.

\subsection*{Answer}

Start with the Yule-Walker equations:
\[
\rho_1 = \phi_{21} + \phi_{22}\rho_1.
\]

Rearrange this to find \(\phi_{21}\):
\[
\phi_{21}
=
\rho_1-\phi_{22}\rho_1
=
\rho_1(1-\phi_{22}).
\]

Now use the second equation:
\[
\rho_2 = \phi_{21}\rho_1 + \phi_{22}.
\]

Substitute
\[
\phi_{21}=\rho_1(1-\phi_{22}).
\]

Then
\[
\rho_2
=
\rho_1(1-\phi_{22})\rho_1+\phi_{22}.
\]

So
\[
\rho_2
=
\rho_1^2(1-\phi_{22})+\phi_{22}.
\]

Expand:
\[
\rho_2
=
\rho_1^2-\rho_1^2\phi_{22}+\phi_{22}.
\]

Group the \(\phi_{22}\) terms:
\[
\rho_2-\rho_1^2
=
\phi_{22}(1-\rho_1^2).
\]

Therefore,
\[
\boxed{
\phi_{22}
=
\frac{\rho_2-\rho_1^2}{1-\rho_1^2}
}.
\]

This is the lag 2 PACF.

\subsection*{Simple Interpretation}

The ACF at lag 2, \(\rho_2\), includes both direct and indirect effects.  The
indirect effect can pass through lag 1.  The PACF removes that lag 1 pathway and
keeps only the direct lag 2 effect.

\section*{Question 1(c): Variance of an AR(2) Process}

\subsection*{Theory to Use}

Consider the AR(2) model
\[
Y_t = \phi_0+\phi_1Y_{t-1}+\phi_2Y_{t-2}+e_t,
\]
where
\[
e_t \sim WN(0,\sigma_e^2).
\]

If the process is stationary, the mean is constant:
\[
\mu = E(Y_t).
\]

Taking expectation:
\[
\mu=\phi_0+\phi_1\mu+\phi_2\mu.
\]

Thus,
\[
\mu = \frac{\phi_0}{1-\phi_1-\phi_2}.
\]

The intercept \(\phi_0\) affects the mean, but it does not affect the variance.

For a stationary AR(2), the variance is
\[
\gamma_0 = \operatorname{Var}(Y_t).
\]

The Yule-Walker equations are
\[
\gamma_1 = \phi_1\gamma_0+\phi_2\gamma_1,
\]
and
\[
\gamma_0 = \phi_1\gamma_1+\phi_2\gamma_2+\sigma_e^2.
\]

Also,
\[
\gamma_2 = \phi_1\gamma_1+\phi_2\gamma_0.
\]

\subsection*{Question}

Suppose
\[
Y_t = \phi_0+\phi_1Y_{t-1}+\phi_2Y_{t-2}+e_t.
\]
Find \(V(Y_t)\).

\subsection*{Answer}

Let
\[
\gamma_0 = V(Y_t).
\]

From the first Yule-Walker equation,
\[
\gamma_1 = \phi_1\gamma_0+\phi_2\gamma_1.
\]

Move the \(\phi_2\gamma_1\) term to the left:
\[
\gamma_1-\phi_2\gamma_1=\phi_1\gamma_0.
\]

So
\[
\gamma_1(1-\phi_2)=\phi_1\gamma_0.
\]

Hence,
\[
\gamma_1=\frac{\phi_1}{1-\phi_2}\gamma_0.
\]

Now use
\[
\gamma_0=\phi_1\gamma_1+\phi_2\gamma_2+\sigma_e^2.
\]

Also,
\[
\gamma_2=\phi_1\gamma_1+\phi_2\gamma_0.
\]

Substitute \(\gamma_2\) into the equation for \(\gamma_0\):
\[
\gamma_0
=
\phi_1\gamma_1+\phi_2(\phi_1\gamma_1+\phi_2\gamma_0)+\sigma_e^2.
\]

Expand:
\[
\gamma_0
=
\phi_1\gamma_1+\phi_1\phi_2\gamma_1+\phi_2^2\gamma_0+\sigma_e^2.
\]

Group the first two terms:
\[
\gamma_0
=
\phi_1(1+\phi_2)\gamma_1+\phi_2^2\gamma_0+\sigma_e^2.
\]

Substitute
\[
\gamma_1=\frac{\phi_1}{1-\phi_2}\gamma_0.
\]

Then
\[
\gamma_0
=
\phi_1(1+\phi_2)
\left(\frac{\phi_1}{1-\phi_2}\gamma_0\right)
+\phi_2^2\gamma_0+\sigma_e^2.
\]

So
\[
\gamma_0
=
\frac{\phi_1^2(1+\phi_2)}{1-\phi_2}\gamma_0
+\phi_2^2\gamma_0+\sigma_e^2.
\]

Move the \(\gamma_0\) terms to the left:
\[
\gamma_0
\left[
1-\phi_2^2-\frac{\phi_1^2(1+\phi_2)}{1-\phi_2}
\right]
=
\sigma_e^2.
\]

Therefore,
\[
\gamma_0
=
\frac{\sigma_e^2}{
1-\phi_2^2-\dfrac{\phi_1^2(1+\phi_2)}{1-\phi_2}
}.
\]

This can be simplified to
\[
\boxed{
V(Y_t)
=
\gamma_0
=
\frac{\sigma_e^2(1-\phi_2)}
{(1+\phi_2)\left((1-\phi_2)^2-\phi_1^2\right)}
}.
\]

Equivalently,
\[
\boxed{
V(Y_t)
=
\frac{\sigma_e^2(1-\phi_2)}
{(1+\phi_2)(1-\phi_2-\phi_1)(1-\phi_2+\phi_1)}
}.
\]

\begin{remark}
This formula is valid when the AR(2) process is stationary.  If the process is
not stationary, a finite constant variance may not exist.
\end{remark}

\subsection*{Simple Interpretation}

The intercept \(\phi_0\) shifts the process up or down.  It changes the mean.
But variance measures spread around the mean, so \(\phi_0\) does not appear in
the final variance formula.

\section*{Question 1(d): Stationarity from the Characteristic Polynomial}

\subsection*{Theory to Use}

For an AR model, stationarity is checked using the characteristic polynomial.
The usual rule is:

\[
\boxed{\text{The AR model is stationary if all roots of the characteristic polynomial lie outside the unit circle.}}
\]

That means every root \(x\) must satisfy
\[
|x|>1.
\]

If any root satisfies
\[
|x|\leq 1,
\]
then the model is not stationary.

\subsection*{Question}

An AR model has characteristic polynomial
\[
(1+0.2x+0.7x^2)(1-0.8x^{12}).
\]

Is the model stationary?

\subsection*{Answer}

We need to check the roots of both factors:
\[
1+0.2x+0.7x^2
\]
and
\[
1-0.8x^{12}.
\]

\subsubsection*{First Factor}

The first factor is
\[
1+0.2x+0.7x^2=0.
\]

Rearrange:
\[
0.7x^2+0.2x+1=0.
\]

For a quadratic
\[
ax^2+bx+c=0,
\]
if the roots are complex conjugates, the product of the roots is
\[
\frac{c}{a}.
\]

Here,
\[
a=0.7,\qquad c=1.
\]

So the product of the two roots is
\[
\frac{1}{0.7}.
\]

Because the roots are complex conjugates, they have the same modulus.  If each
root has modulus \(r\), then
\[
r^2=\frac{1}{0.7}.
\]

Thus,
\[
r=\sqrt{\frac{1}{0.7}}.
\]

Since
\[
\sqrt{\frac{1}{0.7}}>1,
\]
the roots of the first factor lie outside the unit circle.

\subsubsection*{Second Factor}

The second factor is
\[
1-0.8x^{12}=0.
\]

Then
\[
0.8x^{12}=1.
\]

So
\[
x^{12}=\frac{1}{0.8}=1.25.
\]

Therefore, all roots have modulus
\[
|x|=(1.25)^{1/12}.
\]

Since
\[
1.25>1,
\]
we have
\[
(1.25)^{1/12}>1.
\]

So the roots of the second factor also lie outside the unit circle.

\subsubsection*{Conclusion}

All roots of
\[
(1+0.2x+0.7x^2)(1-0.8x^{12})
\]
lie outside the unit circle.

Therefore,
\[
\boxed{\text{the AR model is stationary.}}
\]

\subsection*{Simple Interpretation}

Stationarity means shocks die out over time instead of growing.  The root test
checks whether the AR feedback is weak enough for the process to settle around a
constant mean and variance.  Since every root is outside the unit circle, the
model is stable and stationary.

\section*{Extra Practice Questions}

\subsection*{Practice 1: ACF of an MA(1)}

\textbf{Question.} Suppose
\[
Y_t=e_t+0.6e_{t-1},
\]
where
\[
e_t\sim WN(0,4).
\]
Find the ACF.

\textbf{Answer.}

Here,
\[
\psi=0.6,\qquad \sigma_e^2=4.
\]

For an MA(1),
\[
\gamma_0=(1+\psi^2)\sigma_e^2.
\]

Thus,
\[
\gamma_0=(1+0.6^2)(4)=1.36(4)=5.44.
\]

Also,
\[
\gamma_1=\psi\sigma_e^2=0.6(4)=2.4.
\]

So
\[
\rho_1=\frac{\gamma_1}{\gamma_0}
=
\frac{2.4}{5.44}
=
0.4412.
\]

Therefore,
\[
\boxed{
\rho_k=
\begin{cases}
1, & k=0,\\
0.4412, & |k|=1,\\
0, & |k|\geq 2.
\end{cases}
}
\]

\subsection*{Practice 2: Lag 2 PACF}

\textbf{Question.} If
\[
\rho_1=0.5,\qquad \rho_2=0.2,
\]
find the lag 2 PACF.

\textbf{Answer.}

Use
\[
\phi_{22}=\frac{\rho_2-\rho_1^2}{1-\rho_1^2}.
\]

Substitute:
\[
\phi_{22}
=
\frac{0.2-(0.5)^2}{1-(0.5)^2}.
\]

So
\[
\phi_{22}
=
\frac{0.2-0.25}{1-0.25}
=
\frac{-0.05}{0.75}.
\]

Therefore,
\[
\boxed{\phi_{22}=-0.0667}.
\]

\subsection*{Practice 3: Stationarity Check}

\textbf{Question.} An AR model has characteristic polynomial
\[
1-1.2x.
\]
Is the model stationary?

\textbf{Answer.}

Set the characteristic polynomial equal to zero:
\[
1-1.2x=0.
\]

Then
\[
1.2x=1.
\]

So
\[
x=\frac{1}{1.2}=0.8333.
\]

Since
\[
|x|<1,
\]
the root is inside the unit circle.

Therefore,
\[
\boxed{\text{the model is not stationary.}}
\]

\end{document}
